\documentclass[12pt,a4paper]{article}
\usepackage[utf8]{inputenc}
\usepackage{amsmath}
\usepackage{amsfonts}
\usepackage{amssymb}
\usepackage{graphicx}
\usepackage{geometry}
\usepackage{hyperref}
\usepackage{listings}
\usepackage{color}
\usepackage{booktabs}
\usepackage{fancyhdr}
\usepackage{titlesec}

% Page Geometry
\geometry{a4paper, margin=1in}

% Header and Footer Style
\pagestyle{fancy}
\fancyhf{}
\rhead{\small Dynamic Service Routing and Traffic Management}
\lhead{\small CSE/IT Academic Report}
\rfoot{\thepage}

% Title Formatting
\titleformat{\section}{\normalfont\Large\bfseries}{\thesection.}{1em}{}
\titleformat{\subsection}{\normalfont\large\bfseries}{\thesubsection.}{1em}{}

% Code Listings Configuration
\definecolor{codegray}{rgb}{0.5,0.5,0.5}
\definecolor{codegreen}{rgb}{0,0.5,0}
\definecolor{codeblue}{rgb}{0,0,0.7}
\definecolor{backcolour}{rgb}{0.97,0.97,0.95}

\lstdefinestyle{mystyle}{
    backgroundcolor=\color{backcolour},   
    commentstyle=\color{codegreen},
    keywordstyle=\color{codeblue},
    numberstyle=\tiny\color{codegray},
    stringstyle=\color{magenta},
    basicstyle=\ttfamily\footnotesize,
    breakatwhitespace=false,         
    breaklines=true,                 
    captionpos=b,                    
    keepspaces=true,                 
    numbers=left,                    
    numbersep=5pt,                  
    showspaces=false,                
    showstringspaces=false,
    showtabs=false,                  
    tabsize=2
}
\lstset{style=mystyle}

\begin{document}

% --- Cover Page ---
\begin{titlepage}
    \centering
    \vspace*{1.5cm}
    {\huge\bfseries ACADEMIC PROJECT REPORT \par}
    \vspace{0.8cm}
    {\Large\bfseries AUTOMATED BLUE-GREEN DEPLOYMENT PIPELINE WITH TELEMETRY SCRAPING IN KUBERNETES \par}
    \vspace{2.5cm}
    {\large Submitted in partial fulfillment of the requirements for the course of study \par}
    \vspace{0.3cm}
    {\large\bfseries Bachelor of Technology \par}
    {\large in \par}
    {\large\bfseries Computer Science \& Engineering / Information Technology \par}
    \vspace{3.5cm}
    \begin{minipage}{0.4\textwidth}
        \begin{flushleft} \large
            \emph{Prepared By:}\\
            \textbf{Yasir} \\
            Student
        \end{flushleft}
    \end{minipage}
    \hfill
    \begin{minipage}{0.4\textwidth}
        \begin{flushright} \large
            \emph{Guided By:} \\
            \textbf{Department Faculty} \\
            Department of CSE/IT
        \end{flushright}
    \end{minipage}
    \vfill
    {\large Department of Computer Science \& Engineering \par}
    \vspace{0.5cm}
    {\large \today \par}
\end{titlepage}

\tableofcontents
\newpage

% --- 1. Introduction ---
\section{Introduction}
Modern software release workflows prioritize high availability and minimal service disruption during upgrades. This project describes the implementation of a local continuous integration and continuous delivery (CI/CD) pipeline that automates a Blue-Green deployment strategy. The environment is hosted on a local Kubernetes cluster (Minikube). It consists of a containerized React (Vite) frontend application and an Express (Node.js) backend API. Terraform is used to provision the Kubernetes namespaces. Automation is managed by a Jenkins pipeline that runs automated health checks and controls traffic routing via service selector patches. Application instrumentation is accomplished using a Prometheus client to expose request metrics for collection and visualization via Grafana.

% --- 2. Profile of the Problem ---
\section{Profile of the Problem. Rationale/Scope of the study}
\subsection{Problem Statement}
During typical software updates, standard in-place deployments frequently introduce service interruptions, compatibility discrepancies, or prolonged recovery windows if a newly published application version fails. Specifically, the challenges include:
\begin{enumerate}
    \item \textbf{Transition State Latency}: Taking resources offline during upgrade cycles creates a service gap.
    \item \textbf{High-risk Rollbacks}: Reverting a failed deployment requires redeploying old source code or manually modifying configurations, which increases system recovery time.
    \item \textbf{Lack of Instantaneous Telemetry}: The absence of immediate performance metrics during the transition makes it difficult to detect errors during the update process.
\end{enumerate}

\subsection{Rationale \& Scope of the Study}
This project implements a Blue-Green deployment model to resolve these issues. By maintaining two identical production environments (named Blue and Green), the active environment can be switched using service selector modifications. 

The scope of this project includes:
\begin{itemize}
    \item Deploying isolated replicas of the application (frontend and backend tiers) in a dedicated Kubernetes namespace.
    \item Scripting a Jenkins pipeline to automate building, testing, routing transitions, and fail-safe recovery logic.
    \item Using Prometheus to gather runtime metrics and using Grafana to visualize HTTP transaction rates and error anomalies.
\end{itemize}

% --- 3. Existing System ---
\section{Existing System}
\subsection{Introduction}
Traditional staging-to-production pipelines deploy upgrades directly in-place (Rolling Updates or Recreate Deployments). These models suffer from transitional instabilities, state mismatch, and lagging user sessions due to browser caching.

\subsection{Existing Software}
Standard continuous delivery tools (e.g., raw bash scripting combined with manual virtual machine updates) lack immediate traffic-switching orchestration and feedback loops from runtime monitoring systems.

\subsection{DFD for Present System}
The traditional manual deployment pipeline does not include continuous monitoring feedback loops or isolated candidate verification. This flow is illustrated below:
\begin{center}
    \textit{[Refer to Appendix for Present System DFD code]}
\end{center}

\subsection{What’s new in the system to be developed}
\begin{enumerate}
    \item \textbf{Label-Based Traffic Routing}: Kubernetes Services are patched at runtime to direct internal and external traffic via label selectors.
    \item \textbf{Non-Intrusive Validation}: Newly deployed images are tested in the active cluster prior to complete traffic shifting.
    \item \textbf{Active Cache Defeating}: The web server (Nginx) configuration explicitly includes headers to suppress browser-level caching of static assets.
    \item \textbf{Automated Telemetry}: Real-time response statistics are scraped directly from application endpoints to identify transition anomalies.
\end{enumerate}

% --- 4. Problem Analysis ---
\section{Problem Analysis}
\subsection{Product Definition}
The system is an automated, local Kubernetes-hosted software delivery pipeline. It includes two identical production-like environments, automated integration scripts, and real-time metric scrapers to manage software releases safely.

\subsection{Feasibility Analysis}
\begin{itemize}
    \item \textbf{Technical Feasibility}: The infrastructure is built with standard DevOps tools: Docker for containerization, Kubernetes for pod orchestration, Terraform for workspace isolation, and Jenkins for pipeline management. This setup is technically feasible on general-purpose hardware.
    \item \textbf{Operational Feasibility}: Automated pipeline tasks run in response to version control pushes. System administrators do not need to manually configure infrastructure during releases.
    \item \textbf{Economic Feasibility}: The pipeline uses open-source software, making it economically practical to deploy and manage without license costs.
\end{itemize}

\subsection{Project Plan}
The schedule represents the critical phases of requirement gathering, design, coding, integration, and validation:
\begin{center}
    \textit{[Refer to Appendix for Project Gantt Chart code]}
\end{center}

% --- 5. Software Requirement Analysis ---
\section{Software Requirement Analysis}
\subsection{Introduction}
This section defines the hardware and software resource limits needed to run the local cluster and pipeline components.

\subsection{General Description}
The target application is structured as a two-tier microservice:
\begin{itemize}
    \item \textbf{Frontend Tiers}: A static React build wrapped inside an Nginx instance.
    \item \textbf{Backend Tiers}: A Node.js API server exposing operational endpoints and gathering Prometheus metrics.
\end{itemize}

\subsection{Specific Requirements}
\subsubsection{Software Requirements}
\begin{itemize}
    \item \textbf{Hypervisor/Container Engine}: Docker Engine v20.10+
    \item \textbf{Local Kubernetes Cluster}: Minikube v1.28+ (supporting ingress and metrics-server addons)
    \item \textbf{Infrastructure Provisioner}: Terraform v1.0.0+
    \item \textbf{Automation Platform}: Jenkins LTS (installed locally on the host machine or run as a service)
    \item \textbf{Monitoring Stack}: Prometheus v2.40+ and Grafana v9.0+
\end{itemize}

\subsubsection{Hardware Requirements}
\begin{itemize}
    \item \textbf{Processor}: Multi-core x86-64 Processor (4 physical cores minimum)
    \item \textbf{RAM}: 16 GB system memory (allocated to support Minikube, Docker, and the local Jenkins engine)
    \item \textbf{Disk Space}: 40 GB solid-state drive space (SSD) for caching local container images
\end{itemize}

% --- 6. Design ---
\section{Design}
\subsection{System Design}
The system layout uses label selectors inside Kubernetes Service resources to map incoming ports to either Blue or Green deployment targets dynamically:
\begin{center}
    \textit{[Refer to Appendix for System Design Architecture Diagram code]}
\end{center}

\subsection{Design Notations}
\begin{itemize}
    \item \textbf{Pods}: Groupings of containerized application runtimes.
    \item \textbf{Services}: High-level abstractions defining persistent ingress points mapping to target pods.
    \item \textbf{Namespaces}: Logical isolation domains splitting demo resources from system processes.
\end{itemize}

\subsection{Detailed Design}
\subsubsection{Static File Routing and Caching Overrides}
The frontend reverse proxy is configured via an Nginx block that forwards API queries to the backend service. It injects specific control headers to bypass intermediate cache servers.
\subsubsection{Prometheus Scraping Target}
An active scrape job is declared in the monitoring namespace, pointing to the internal DNS name of the backend: \texttt{backend-service.blue-green-demo.svc.cluster.local:5000/metrics}.

\subsection{Flowcharts}
The pipeline flowchart tracks the execution path from git trigger down to deployment checks and rollbacks:
\begin{center}
    \textit{[Refer to Appendix for Jenkins Pipeline Flowchart code]}
\end{center}

\subsection{Pseudo Code}
The Jenkinsfile pipeline stages are modeled in the procedural block below:
\begin{lstlisting}[language=Python, caption=Jenkins Pipeline Control Loop]
pipeline_execution():
    # Stage 1: Setup Namespace via Terraform
    directory_change("terraform")
    terraform_init()
    terraform_apply()
    directory_change("..")
    
    # Stage 2: Verify and Deploy Blue Base
    docker_env_eval()
    if not kubernetes_deployment_exists(name="blue-frontend", namespace="blue-green-demo"):
        build_docker_image(tag="devops-proj/frontend:blue", build_arg="VERSION=blue", path="./frontend")
        build_docker_image(tag="devops-proj/backend:blue", build_arg="VERSION=blue", path="./backend")
        kubernetes_apply("k8s/blue-deployment.yaml")
        kubernetes_apply("k8s/service.yaml")
        thread_sleep(10)
        
    # Stage 3: Build Green Images
    try:
        build_docker_image(tag="devops-proj/frontend:green", build_arg="VERSION=green", path="./frontend")
        build_docker_image(tag="devops-proj/backend:green", build_arg="VERSION=green", path="./backend")
    except BuildException as e:
        pipeline_abort(message="Green Image Build Failed. System Unchanged.")
        
    # Stage 4: Deploy Green Candidate
    kubernetes_apply("k8s/green-deployment.yaml")
    
    # Stage 5: Shift Routing Selector to Green
    kubernetes_patch_service(service="frontend-service", selector="green")
    kubernetes_patch_service(service="backend-service", selector="green")
    
    # Stage 6: Perform Integration Health Check
    thread_sleep(12)
    minikube_host = get_minikube_ip()
    http_response_code = execute_http_get(url=f"http://{minikube_host}:30080/api/health")
    
    # Stage 7: Evaluate Validation Response
    if http_response_code != 200:
        log_warning("Health Check Failed! Restoring stable Blue state...")
        kubernetes_patch_service(service="frontend-service", selector="blue")
        kubernetes_patch_service(service="backend-service", selector="blue")
        kubernetes_delete_deployment(name=["green-frontend", "green-backend"])
        terminate_pipeline(status="FAILED")
    else:
        log_info("Deployment Successful. Green version promoted to active production.")
        terminate_pipeline(status="SUCCESS")
\end{lstlisting}

% --- 7. Testing ---
\section{Testing}
\subsection{Functional Testing}
Tests verified routing behavior by polling the application version endpoint. The target returned the assigned environment label (`blue` or `green`) based on the current service routing configuration.

\subsection{Structural Testing}
To test the pipeline's recovery behavior, a fault was injected into the backend application source code:
\begin{lstlisting}[language=JavaScript]
// Backend simulated failure
app.get('/api/health', (req, res) => {
  return res.status(500).json({ status: 'error', version: VERSION });
});
\end{lstlisting}
The pipeline successfully detected the non-200 HTTP code during execution and automatically rolled back traffic to the stable Blue pods.

\subsection{Levels of Testing}
\begin{itemize}
    \item \textbf{Unit Level}: Checked the syntax of components during static builds.
    \item \textbf{Integration Level}: Verified the communication path between frontend proxies and API backend services inside Kubernetes.
    \item \textbf{System Level}: Validated the complete delivery cycle, including triggers, building, and rollback processes.
\end{itemize}

\subsection{Testing the Project}
\begin{table}[h]
\centering
\caption{System Test Matrix}
\begin{tabular}{@{}lllll@{}}
\toprule
\textbf{Test ID} & \textbf{Scenario} & \textbf{Expected Output} & \textbf{Observed Output} & \textbf{Result} \\ \midrule
TC-01 & Healthy Update & Green environment promoted & Traffic routed to Green & PASS \\
TC-02 & Fault Injection & Pipeline triggers rollback & Traffic restored to Blue & PASS \\
TC-03 & Telemetry Verify & Metrics scraped on endpoint & Grafana shows HTTP metrics & PASS \\ \bottomrule
\end{tabular}
\end{table}

% --- 8. Implementation ---
\section{Implementation}
\subsection{Implementation of the project}
The code was deployed using Docker environment hooks inside Minikube. Initial configurations:
\begin{itemize}
    \item \textbf{IaC Configuration}: Terraform files configured inside the \texttt{terraform} directory managed namespace allocations.
    \item \textbf{Manifest Declarations}: Declarative files in the \texttt{k8s} folder mapped deployment specs and NodePorts.
    \item \textbf{Alerting Config}: Custom rules configured in the Prometheus stack logged network rates and tracked anomalies.
\end{itemize}

\subsection{Conversion Plan}
Upgrading the infrastructure to use Blue-Green deployments was handled via a **parallel transition strategy**:
\begin{itemize}
    \item The active Blue deployment handles client requests.
    \item The Green environment is deployed and verified in isolation within the same namespace.
    \item A patch command updates service selectors to swap active targets.
\end{itemize}

\subsection{Post-Implementation and Software Maintenance}
Declarative files are tracked in source control. Monitoring parameters inside Grafana are reviewed regularly to verify node utilization limits.

% --- 9. Project Legacy ---
\section{Project Legacy}
\subsection{Current Status of the project}
The pipeline is fully operational locally. Build schedules are triggered on push events via ngrok tunnels forwarding webhooks directly from GitHub repositories to the local Jenkins daemon.

\subsection{Remaining Areas of concern}
\begin{itemize}
    \item \textbf{Persistent State Databases}: High-availability database replication configurations are needed to prevent data collision between blue/green versions.
    \item \textbf{Auto-Scaling Rules}: Adding Horizontal Pod Autoscalers (HPA) to scale pod replicas based on traffic load dynamically.
\end{itemize}

\subsection{Technical and Managerial lessons learnt}
\begin{itemize}
    \item \textbf{Fail-Fast Principles}: Catching application failures early in the CI stage protects users from system instability.
    \item \textbf{Infrastructure as Code (IaC)}: Declarative configurations guarantee environment stability, preventing discrepancies between staging and production stages.
\end{itemize}

% --- 10. User Manual ---
\section{User Manual}
\subsection{Step 1: Initialize Cluster}
\begin{lstlisting}[language=bash]
minikube start --driver=docker
minikube addons enable ingress
minikube addons enable metrics-server
eval $(minikube docker-env)
\end{lstlisting}

\subsection{Step 2: Deploy Telemetry Stack}
\begin{lstlisting}[language=bash]
helm repo add prometheus-community https://prometheus-community.github.io/helm-charts
helm repo update
helm install monitoring prometheus-community/kube-prometheus-stack --namespace monitoring --create-namespace
kubectl apply -f k8s/servicemonitor.yaml
kubectl apply -f prometheus/alert-rules.yaml
\end{lstlisting}

\subsection{Step 3: Expose dashboards}
\begin{lstlisting}[language=bash]
kubectl port-forward svc/monitoring-grafana 3000:80 -n monitoring
# Grafana user: admin | pass: prom-operator
\end{lstlisting}

% --- 11. Source Code ---
\section{Source Code (System Snapshots)}
\subsection{Dockerfile (Conditional Build logic)}
\begin{lstlisting}[language=HTML, caption=Dockerfile]
FROM node:18-alpine AS build
WORKDIR /app
COPY package*.json ./
RUN npm install
COPY . .

ARG VERSION=blue
RUN if [ "$VERSION" = "green" ]; then cp src-green/App.jsx src/App.jsx; fi

RUN npm run build
FROM nginx:alpine
COPY --from=build /app/dist /usr/share/nginx/html
COPY nginx.conf /etc/nginx/conf.d/default.conf
EXPOSE 80
CMD ["nginx", "-g", "daemon off;"]
\end{lstlisting}

\subsection{Express Telemetry setup (\texttt{backend/app.js})}
\begin{lstlisting}[language=JavaScript, caption=Express Prometheus tracking]
const client = require('prom-client');
const register = new client.Registry();
client.collectDefaultMetrics({ register });

const httpRequestsTotal = new client.Counter({
  name: 'http_requests_total',
  help: 'Total number of HTTP requests processed',
  labelNames: ['method', 'handler', 'status']
});
register.registerMetric(httpRequestsTotal);
\end{lstlisting}

% --- 12. Bibliography ---
\section{Bibliography}
\begin{enumerate}
    \item \textbf{Kubernetes Up \& Running}: Dive into the Future of Infrastructure, 2nd Edition - Brendan Burns, Joe Beda, Kelsey Hightower.
    \item \textbf{Terraform: Up \& Running}: Writing Infrastructure as Code - Yevgeniy Brikman.
    \item \textbf{Continuous Delivery}: Reliable Software Releases through Build, Test, and Deployment Automation - Jez Humble, David Farley.
    \item \textbf{Prometheus Documentation}: Monitoring Systems and Alerting Toolkits (https://prometheus.io/docs/).
\end{enumerate}

\newpage
% --- Appendix: Mermaid Diagrams Source Code ---
\section*{Appendix: Mermaid Diagrams Source Code}
\addcontentsline{toc}{section}{Appendix: Mermaid Diagrams Source Code}

\subsection*{Diagram 1: Data Flow Diagram (DFD) of Existing System}
\begin{lstlisting}
graph TD
    Developer[Developer] -->|Manual Git Push| GitRepo[Git Repository]
    GitRepo -->|Trigger webhook| OldCI[Basic CI Server]
    OldCI -->|Direct in-place build| ProductionServer[Production Server]
    ProductionServer -->|Errors occur: System crashes| Users[Users Experience Downtime]
    Users -->|Manual Complaint| Developer
    Developer -->|Manual Revert & Restart| ProductionServer
\end{lstlisting}

\subsection*{Diagram 2: Gantt Chart Project Plan}
\begin{lstlisting}
gantt
    title Blue-Green Deployment System Project Schedule
    dateFormat  YYYY-MM-DD
    section Requirement Analysis & Research
    System Study & IaC Selection        :active, 2026-05-01, 3d
    section Architecture & Design
    Design Kubernetes YAMLs & Services :2026-05-04, 4d
    Jenkins Pipeline Definition        :2026-05-08, 4d
    section Implementation
    Build Frontend & Backend           :2026-05-12, 5d
    Terraform Code & Namespace setup   :2026-05-17, 2d
    Monitoring Integration             :2026-05-19, 3d
    section Testing & Deployment
    Pipeline Dry Runs & Fail-safes     :2026-05-22, 2d
\end{lstlisting}

\subsection*{Diagram 3: System Design Architecture}
\begin{lstlisting}
graph TD
    User([End User]) -->|Accesses NodePort:30080| FrontendService[Kubernetes frontend-service]
    FrontendService -->|Selectors label: version=blue/green| ActiveDeployment{Active Selector}
    ActiveDeployment -->|version=blue| BlueFrontend[blue-frontend Pod]
    ActiveDeployment -->|version=green| GreenFrontend[green-frontend Pod]
    
    BlueFrontend -->|Proxy /api| BackendService[Kubernetes backend-service]
    GreenFrontend -->|Proxy /api| BackendService
    
    BackendService -->|version=blue/green| ActiveBackend{Active Selector}
    ActiveBackend -->|version=blue| BlueBackend[blue-backend Pod]
    ActiveBackend -->|version=green| GreenBackend[green-backend Pod]

    Prometheus[Prometheus Server] -->|Scrapes /metrics| BlueBackend
    Prometheus -->|Scrapes /metrics| GreenBackend
    Grafana[Grafana Dashboard] -->|Visualizes Metrics| Prometheus
\end{lstlisting}

\subsection*{Diagram 4: Jenkins Pipeline Flowchart}
\begin{lstlisting}
flowchart TD
    Start([Git Push Triggered]) --> InitTF[Terraform provisions Namespace]
    InitTF --> BuildBlue{Blue Deployment Running?}
    BuildBlue -- No --> DeployBlue[Build & Deploy Blue Version v1.0]
    BuildBlue -- Yes --> BuildGreen[Build Green Image v2.0]
    DeployBlue --> BuildGreen
    BuildGreen --> DeployGreen[Deploy Green Pods]
    DeployGreen --> SwitchTraffic[Switch Service Selector to 'green']
    SwitchTraffic --> Stabilization[Wait for 12s Stabilization]
    Stabilization --> HealthCheck{HTTP GET /api/health == 200?}
    HealthCheck -- Yes (200 OK) --> Complete[Pipeline Finished - Traffic remains on Green]
    HealthCheck -- No (5xx Error) --> Rollback[Switch Service Selector back to 'blue']
    Rollback --> DeleteGreen[Terminate Green Pods]
    DeleteGreen --> Fail([Build Status: Failed - Restored Stable Blue])
\end{lstlisting}

\end{document}
